% !PRINTED 585
\weberpage{585}

Wir wollen auch hier den in der Formel (34) ausgedrückten
Satz hervorheben:

\gesperrt{Jede Primzahl von der Form} $4f+1$ \gesperrt{lässt sich in
die Summe zweier Quadrate zerlegen.}

\begin{center}
\S.~173.\\[6pt]
\gesperrt{Die complexen Zahlen von Gauss.}
\end{center}

Der zuletzt bewiesene Satz bildet die Grundlage für die
Theorie des Zahlkörpers, der durch Adjunction der imaginären
Einheit $i=\sqrt{-1}$ aus dem Körper der rationalen Zahlen entsteht,
den wir nach unserer Festsetzung mit $R(i)$ zu bezeichnen
haben\textsuperscript{1}).

Wir haben hier die Primzahlen von der Form $4f+1$ zu
unterscheiden von denen der Form $4f+3$ und wir wollen der
Kürze wegen die ersten mit der Primzahl 2 zusammenfassen und
mit $p$, die zweiten mit $q$ bezeichnen, also
\[
p=4f+1,2;\quad q=4f+3
\]
setzen. Es gelten dann folgende Sätze:

1. \gesperrt{Für jedes $p$ giebt es zwei ganze Zahlen $a$ und $b$,
so dass}
\[
p=a^2+b^2
\]
\gesperrt{ist.}

Dies ist für die Primzahlen von der Form $4f+1$ im Schlusssatze
 des letzten Paragraphen bewiesen und ist für $p=2$ aus
$2=1^2+1^2$ unmittelbar ersichtlich.

Dem steht ein zweiter Satz gegenüber:

2. \gesperrt{Eine Primzahl $q$ ist niemals in der Form $a^2+b^2$
darstellbar,} oder der noch allgemeinere:

3. \gesperrt{Die Summe zweier Quadrate ganzer Zahlen $a^2+b^2$
ist nur dann durch eine Primzahl $q$ theilbar, wenn $a$ und $b$
durch $q$ theilbar sind.}

Der Satz 2. ergiebt sich einfach daraus, dass, wenn $a^2+b^2$
ungerade ist, die eine der beiden Zahlen $a$, $b$ gerade, also ihr
Quadrat durch 4 theilbar, die andere ungerade, also ihr Quadrat
%CONT>>

\textsuperscript{1}) \gesperrt{Gauss}, Theoria residuorum biquadraticorum, commentatio
secunda. Werke, Bd.~II.
